\documentclass[10pt,a4paper]{article}

\usepackage[T1]{fontenc}
\usepackage[utf8]{inputenc}
\usepackage{lmodern}
\usepackage{amsmath,amssymb,amsthm}
\usepackage{booktabs,array}
\usepackage{geometry}
\usepackage{xcolor}
\usepackage{enumitem}
\usepackage{microtype}
\usepackage{hyperref}
\usepackage{fancyhdr}

\geometry{margin=1in}
\hypersetup{
  colorlinks=true,
  linkcolor=blue!55!black,
  citecolor=blue!55!black,
  urlcolor=blue!55!black,
  pdftitle={Full Distance, Short Detour: A One-Query Forgery against FDSD},
  pdfauthor={Martin Feussner}
}
\setlist{nosep,leftmargin=*}
\setlength{\parindent}{1em}
\setlength{\parskip}{0.25em}
\pagestyle{fancy}
\fancyhf{}
\fancyhead[L]{Full Distance, Short Detour}
\fancyhead[R]{14 October 2026}
\fancyfoot[C]{\thepage}
\renewcommand{\headrulewidth}{0.3pt}

\newtheorem{theorem}{Theorem}
\newtheorem{lemma}{Lemma}
\newcommand{\F}{\mathbb{F}}
\newcommand{\Col}{\mathsf{Col}}
\newcommand{\wt}{\mathsf{wt}}

\title{\bfseries Full Distance, Short Detour: A One-Query Forgery against FDSD}
\author{Martin Feussner\\
\small Selmer Center, University of Bergen\\
\small \texttt{martin.feussner@uib.no}}
\date{14 October 2026}

\begin{document}
\maketitle

\noindent\fcolorbox{black!50}{black!4}{\parbox{0.94\textwidth}{
This audit used the OpenAI model Daybreak Blue at the ultra reasoning level.
OpenAI Codex generated and selected this paper's title and intentionally chose
``Full Distance, Short Detour'' as wordplay contrasting the scheme's
full-distance decoding premise with the short public linear detour used by the
forgery. OpenAI Codex performed AI-assisted target selection, reconstruction,
cryptanalysis, implementation, experiments, and technical writing. A separate
AI agent performed adversarial review on the retained key, and a fresh-context
AI agent independently reproduced the attack on five fresh keys. The final
exact-target history recheck and independent Publication Reviewer gate also
passed. All required publication gates passed, and the result is classified as
\texttt{VERIFIED\_ATTACK}. No independent human verification was performed or
required.}}

\begin{abstract}
We describe a public-only, one-signing-query fresh-message forgery against the
Level-I parameters of Abdelhaliem Babiker's \emph{A Signature Scheme from
Full-Distance Syndrome Decoding}, IACR ePrint 2023/1956. The ordinary verifier
accepts a weight-343 vector $\sigma$ when
\[
  H^\star\sigma=Rh+R_{[25]}\zeta.
\]
The public matrix $H^\star$ retains 125 hidden-support columns spanning
$\Col(R)$. Testing its 1,600 columns for membership in that public subspace
therefore recovers the support and a public map $U$ satisfying
$H^\star U=R$. From one ordinary signature on $h_0$, an attacker transfers the
syndrome equation to any fresh target $h_1$. A 25-bit affine correction leaves
only the verifier's exact-weight condition; reflected Gray-code search finds a
weight-343 representative.

The primary full-parameter run recovered all 125 columns and forged its
source-frozen unsigned target after 3,657 candidates in 0.627 seconds of
internal attack time, using one worker and 19,456 KiB peak resident memory.
An independently written adversarial implementation forged eight additional
frozen targets on the same retained key, and an exhaustive $2^{25}$ sweep for
one target found 1,337 accepted exact-weight representatives. A fresh-context
implementation then succeeded on five of five wholly fresh keys with one
ordinary query per key; the slowest attack took 13.376 seconds and all controls
rejected. A final scheme-specific recheck located no previous public attack
against the exact target. No author code or test vectors were located; all
experiments use an independently documented mathematical reconstruction of
the paper. The technical, practical, adversarial, independent-reproduction,
final-history, and independent-publication-review gates have all passed. The
result is classified as \texttt{VERIFIED\_ATTACK}.
\end{abstract}

\section{Introduction}

The Full-Distance Syndrome Decoding (FDSD) signature proposal presents three
parameter sets intended for 128-, 192-, and 256-bit security~\cite{fdsd}. Its
verification equation resembles a random syndrome-decoding instance: the
signer returns a fixed-weight binary vector, and the verifier checks its
syndrome against a message-dependent public target. The Level-I row is
\[
  (n,k,L,\epsilon,\theta,d)=(1600,400,125,5,25,343).
\]
The target argues that a forger must solve a full-distance syndrome-decoding
problem and uses the Syndrome Decoding Estimator as an indicative parameter
benchmark~\cite{esserbellini}.

The public matrix is not an unstructured decoding instance. Key generation
preserves a 125-coordinate subset on which $H^\star$ has exactly the public
column space of $R$. This exposes both the hidden support and an equivalent
public linear transformation. A queried signature can then be moved to a new
message syndrome. The only nonlinear obstacle is the prescribed weight, and
the verifier gives 25 freely adjustable bits with which to restore it.

The resulting attack has the following scope:
\begin{itemize}
  \item it targets the complete advertised Level-I row;
  \item it uses one ordinary signature on a message distinct from the target;
  \item it reads only the public matrices and that one signature;
  \item it outputs a weight-343 pair accepted by the ordinary verifier on a
        fresh attacker-fixed message; and
  \item it exploits the specified key structure and verification equation,
        rather than a parser, fault, weak random generator, or side channel.
\end{itemize}

\paragraph{Technique attribution.}
We make no novelty claim for Gaussian elimination, public subspace-membership
tests, affine signature transformation, or Gray-code enumeration. The
contribution is the concrete observation and full-parameter demonstration
that FDSD's structured public key exposes the transformation needed to apply
those standard tools and bypass the advertised decoding problem.

\section{Target, reconstruction, and claim boundary}

\subsection{Pinned target}

The target is Abdelhaliem Babiker's \emph{A Signature Scheme from
Full-Distance Syndrome Decoding}, IACR ePrint 2023/1956~\cite{fdsd}. The audit
uses the sole object observed with archive datestamp 24 December 2023. Its
26-page PDF has SHA-256 digest
\begin{center}
\small\ttfamily
394b7a3e4f8d935cfabffd90e17b4c8cdd5708267a14ede67b39c49ab9ead822.
\end{center}
The extracted text has SHA-256 digest
\begin{center}
\small\ttfamily
8763cd064947c50af81e487a9f7abf294b86893d2596cdff94bb199db33e460a.
\end{center}

No author implementation, source repository, byte format, or test vector was
located. The paper is therefore authoritative. The audit built an independent
paper-level adapter and froze the following otherwise unstated conventions:
the balanced group weights are eighteen 14s followed by seven 13s;
$\mathsf{Hash}_{125}$ is the first 125 SHAKE256 output bits in MSB-first byte
order; and mathematical vector coordinates are serialized MSB-first with zero
padding. The attack depends on the paper's algebra, not on these particular
byte choices: it accepts the same $h$ that the verifier receives and operates
entirely over $\F_2$ afterward.

The faithful adapter checked all 125 basic and 800 advanced key identities.
On one fresh key, 4,096 of 4,096 honest signatures had exact weight 343 and
were accepted; changed-message, changed-$\sigma$, changed-$\zeta$, and
changed-public-key controls rejected. The public attack process received only
$H^\star$, $R$, and one exported ordinary signature.

\subsection{Level-I public relation}

All vectors and matrices below are over $\F_2$. Let
\[
 H^\star\in\F_2^{1200\times1600},\qquad
 R\in\F_2^{1200\times125}.
\]
For a message $m$, the reconstructed adapter computes
$h=\mathsf{Hash}_{125}(m)$. The ordinary verifier accepts
$(\sigma,\zeta)\in\F_2^{1600}\times\F_2^{25}$ precisely when
\begin{equation}
  \wt(\sigma)=343,
  \qquad
  H^\star\sigma=Rh+R_{[25]}\zeta,
  \label{eq:verify}
\end{equation}
where $R_{[25]}$ denotes the first 25 columns of $R$.

The claim concerns this mathematical KeyGen/Sign/Verify relation. It does not
assert conformance to absent author software or an absent external signature
encoding. It does not claim recovery of the nominal private key.

\subsection{Attack-history status}

We located no previous public attack against Abdelhaliem Babiker's exact
ePrint 2023/1956 scheme in the searches completed on 14 October 2026. The
search covered exact and normalized titles, identifier, author, aliases,
citation indexes, public code and issue surfaces, related work, the distinct
2022/1698 predecessor, and the specific public-column recovery route. This is
a bounded negative result: it cannot exclude private, deleted, or unindexed
work. OpenAlex did not index the target; Semantic Scholar returned HTTP 429;
unauthenticated GitHub code search returned HTTP 401; and the official
all-versions endpoint returned HTTP 403.

The official landing page still lists receipt on 24 December 2023 and approval
on 25 December 2023. A current mirror PDF downloaded on 14 October 2026 was
byte-for-byte identical to the frozen target, with the SHA-256 digest given
above; we located no later target revision. The sealed final-history handoff
has SHA-256
\begin{center}
\small\ttfamily
623f30f514d8095d558b27475dd6fe8cc122b41f84f455f02bbc79f5707d150c,
\end{center}
and the sealed prior-art manifest has SHA-256
\begin{center}
\small\ttfamily
5ffc63c382d18d2d3437092c8a0f637b256d1269f8f3120455334cc5947f0a25.
\end{center}

\section{Public recovery of the hidden signing map}

The basic construction defines a 125-coordinate set $T_e$ and a matrix $U$
such that
\begin{equation}
  HU=R.
  \label{eq:basic}
\end{equation}
The restriction $U_{T_e}$ is a row/column permutation of the nonsingular
middle block $C_1$, so it is invertible. Consequently,
\begin{equation}
  H_{[:,T_e]}U_{T_e}=R,
  \qquad
  \Col(H_{[:,T_e]})=\Col(R).
  \label{eq:span}
\end{equation}

Advanced key generation changes columns indexed outside $T_e$ and explicitly
retains $H^\star_{[:,j]}=H_{[:,j]}$ for $j\in T_e$. It follows that every one
of these 125 public columns lies in $\Col(R)$. The attack computes a basis for
$R$, reduces all 1,600 columns of $H^\star$ against it, and selects the
members. On the retained key it obtains exactly 125 members of rank 125; all
other 1,475 columns lie outside the subspace. For random-looking outside
columns, accidental membership in a 125-dimensional subspace of
$\F_2^{1200}$ has probability $2^{-1075}$ per column. The implementation
nevertheless checks the exact count and rank and aborts on ambiguity.

Let $J$ be the recovered coordinate set. Solving
\[
  H^\star_{[:,J]}x_j=R_j
\]
for every $0\le j<125$, and embedding each $x_j$ into the 1,600 coordinates,
gives public columns $u_j$ with
\begin{equation}
  H^\star u_j=R_j.
  \label{eq:recovered}
\end{equation}
Writing $U=(u_0,\ldots,u_{124})$ yields the public equivalent transformation
$H^\star U=R$. This recovers only the capability required to translate
syndromes; it need not reconstruct the nominal secret-key representation.

\section{One-query fresh-message forgery}

Suppose an ordinary signing query on a non-target message returns
$(\sigma_0,\zeta_0)$ for hash $h_0$. Then
\begin{equation}
  H^\star\sigma_0=Rh_0+R_{[25]}\zeta_0.
\end{equation}
Freeze a fresh attacker-chosen target message with hash $h_1$. Split each
125-bit hash into a 25-bit head and a 100-bit tail. For arbitrary
$a\in\F_2^{25}$ define
\begin{align}
  \delta_{[25]} &= a,\\
  \delta_{[25,125)} &= h_{0,[25,125)}+h_{1,[25,125)},\\
  \sigma_1(a) &= \sigma_0+U\delta,\\
  \zeta_1(a) &= \zeta_0+h_{0,[25]}+h_{1,[25]}+a.
  \label{eq:transform}
\end{align}

\begin{theorem}
For every $a\in\F_2^{25}$, the pair in
Equation~\eqref{eq:transform} satisfies the syndrome equality in
Equation~\eqref{eq:verify} for target hash $h_1$.
\end{theorem}

\begin{proof}
Using $H^\star U=R$,
\[
H^\star\sigma_1
=Rh_0+R_{[25]}\zeta_0+R\delta.
\]
The two copies of $h_{0,[25,125)}$ cancel, while the head contribution is
$R_{[25]}a$. Regrouping the remaining terms gives
\[
H^\star\sigma_1
=Rh_1+R_{[25]}
 (\zeta_0+h_{0,[25]}+h_{1,[25]}+a)
=Rh_1+R_{[25]}\zeta_1.
\]
\end{proof}

Thus the public equation holds identically and the attacker must only find an
$a$ for which $\wt(\sigma_1(a))=343$. The implementation walks the at most
$2^{25}=33{,}554{,}432$ values in reflected Gray-code order. Consecutive
candidates differ by one recovered $u_j$, so each update is one XOR and one
weight test. No syndrome decoding is performed.

\section{Experimental evidence}

\subsection{Primary retained key}

The sealed primary attack fixed its target before reading any public attack
input. It then opened only four files: public metadata, $H^\star$, $R$, and
exactly the first 609-byte JSONL signature record. An interpreter audit hook
enforced that allowlist. No private key, seed, coordinate labels, validation
summary, second signature, or network input was available.

On the fresh retained Level-I key, $R$ had rank 125; exactly 125 columns of
$H^\star$ lay in $\Col(R)$ and had rank 125; and all 125 recovered identities
$H^\star u_j=R_j$ passed. The first exact-weight target representative appeared
after 3,657 candidates. The forged vector had weight 343 and the ordinary
verifier accepted. Reusing the queried signature on the target, changing the
target message, changing $\sigma$, and changing $\zeta$ all rejected.

The attack reported 0.0536 seconds for public recovery, 0.00244 seconds for
the correction search, and 0.627 seconds total internal time. GNU
\texttt{time} measured 1.70 seconds wall time and 19,456 KiB peak RSS on one
worker. A deterministic replay matched every non-timing result field.

\subsection{Adversarial retained-key review}

A separate AI reviewer re-derived and reimplemented the attack with the
opposite pivot convention and independent row- and column-oriented verifiers.
Using the same sole query, it accepted forgeries for eight source-frozen,
previously unsigned targets. First solutions appeared after 6,351 to 88,507
candidates. Changed-message, unchanged-query, one-bit-$\zeta$, and
weight-preserving two-coordinate $\sigma$ controls rejected for all eight.

For one target, the review exhausted all $2^{25}$ corrections and found 1,337
exact-weight representatives, all satisfying the public equation. The random
affine-space heuristic predicts 1,339.56 for the retained query's 41 selected
support coordinates. The full sweep and eight-target review took 23.70
seconds wall time and 19,456 KiB peak RSS on one worker.

\subsection{Independent cross-key reproduction}

A fresh-context AI agent independently implemented the attack without opening
the primary attack lane before its package was sealed. It generated five
wholly fresh Level-I keys from separate \texttt{os.urandom(32)} roots. For
each key it fixed a distinct target before the sole ordinary signing query and
gave a separate attack process only the public key and that one transcript.
All five attacks succeeded, with zero target-message queries.

Every key had rank 125 for $R$, the recovered $H^\star$ restriction, and the
recovered restriction of $U$; every key exposed exactly 125 public columns in
$\Col(R)$. Candidate counts were 303, 1,014, 3,865, 199,806, and 19,856,417
(minimum/median/maximum 303/3,865/19,856,417). Per-key attack totals were
0.828834402, 13.376351877, 0.942027737, 0.835081668, and 0.849683596 seconds.
Each attack process used at most 18,944 KiB RSS. The sequential wrapper,
including five fresh key generations, validation, controls, and sealing, took
41.85 seconds wall time and 22,016 KiB peak RSS on one worker.

For every key the ordinary verifier accepted the weight-343 forgery and the
original query, while changed-message, changed-$\sigma$, changed-$\zeta$,
wrong-key, and unchanged-query-on-target controls rejected. The sealed handoff
has SHA-256
\begin{center}
\small\ttfamily
2c7e621eadd38b5cc1386a64d2398204ac6a7a2afa14f9f7b7933bf3146b0ced,
\end{center}
and the independent package manifest has SHA-256
\begin{center}
\small\ttfamily
e6b391d30065fc8958f58514897814caef9a3658e66eba7934834f674d6f130e.
\end{center}

The public reproducer has also passed an engineering smoke test:
it generated a new full Level-I key without persisting the seed or secret,
exported one ordinary signature, ran the isolated public forger, accepted the
fresh-message forgery under ordinary Verify, and rejected all controls. This
validates package operation independently of the sealed research artifacts.

\section{Complexity and applicability}

Public recovery performs elimination on 125 public 1,200-bit columns, reduces
1,600 public columns, and solves 125 public systems. The final enumeration is
bounded by $2^{25}$ candidates and stores only the public matrices, a basis,
125 recovered columns, and current vectors. The retained runs used about
19 MiB peak RSS and one worker.

The syndrome transformation is deterministic once a queried signature and
target are fixed. Search success depends on the 25-dimensional affine space
containing a representative with the query's required support weight. The
retained key exhibited many solutions and the five independent keys all
yielded a solution, with first hits spanning 303 to 19,856,417 candidates.
These six keys provide empirical evidence rather than a population bound. The
attack detects rank loss or extra subspace columns and aborts rather than
silently misclassifying them.

The route attacks the design structure. It does not depend on the chosen
SHAKE256 adapter beyond obtaining two distinct 125-bit hashes, and it uses no
malformed encoding, implementation error, concurrency, fault, leakage, or
secret state. It presently demonstrates the advertised Level-I row; no claim
is made here for the Level-III or Level-V rows.

\section{Reproducibility}

The public package uses Python~3.11 or later and only the standard library.
Running \path{./reproducer/reproduce.sh} performs these steps in separate
processes:
\begin{enumerate}
  \item generate a fresh complete Level-I key from \texttt{os.urandom(32)},
        sign one non-target message, and persist only public inputs;
  \item freeze the target in the public forger, recover the 125-column map,
        search the 25-bit correction space, and write the forgery; and
  \item load only the public key, query, and result into the ordinary verifier,
        then check freshness and four negative controls.
\end{enumerate}
The successful run ends with the explicit banner
\texttt{ATTACK REPRODUCTION: PASS}. Common numerical-library thread variables
are pinned to one even though no numerical package is used. The generator
never writes its seed, secret key, or coordinate labels. The public package is
covered by a SHA-256 manifest.

\section{Limitations and status}

The evidence establishes a practical, full-parameter, design-level
fresh-message forgery on the retained Level-I key, an eight-target adversarial
review, and a five-of-five fresh-key independent reproduction. The final
exact-target history recheck and independent Publication Reviewer gate also
passed. All required gates passed, and the result is classified as
\texttt{VERIFIED\_ATTACK}.
No independent human verification is required by the project Charter. The
paper makes no universal all-key success-probability claim, no claim beyond
Level I, no claim of nominal secret-key recovery, and no claim that the
general linear-algebra techniques are new.

\section{Conclusion}

FDSD's advanced public key preserves a hidden 125-column basis for the public
space $\Col(R)$. Public subspace membership exposes that basis and reconstructs
a transformation $U$ satisfying $H^\star U=R$. One ordinary signature can then
be translated to any fresh target syndrome, with 25 free bits left to restore
the verifier's exact weight. On the retained full Level-I key the attack took
well under a second internally, a separate adversarial implementation forged
eight further targets, and an independent fresh-context implementation
succeeded on five of five new keys. The advertised full-distance decoding
problem is bypassed by a short public detour through the key's preserved
column space.

\section*{Acknowledgements}
The computations were performed on the Norwegian Research and Education Cloud (NREC), using resources provided by the University of Bergen and the University of Oslo.

\begin{thebibliography}{9}

\bibitem{fdsd}
A.~Babiker.
\newblock A Signature Scheme from Full-Distance Syndrome Decoding.
\newblock IACR Cryptology ePrint Archive, Report 2023/1956, 2023.
\newblock \url{https://eprint.iacr.org/2023/1956}.

\bibitem{esserbellini}
A.~Esser and E.~Bellini.
\newblock Syndrome Decoding Estimator.
\newblock In \emph{Public-Key Cryptography---PKC 2022}, LNCS 13177,
pp.~112--141. Springer, 2022.
\newblock \url{https://doi.org/10.1007/978-3-030-97131-1_5}.

\end{thebibliography}

\end{document}
